% =============================================================================
% Figure 1E-2: AD9742 current-mode DAC operation (simplified)
% =============================================================================
% Conceptual view of the current-source architecture inside the AD9742:
% segmented PMOS current sources sum to a total I_FS set by the FSADJ
% resistor, the 12-bit input code routes each segment to either IOUTA or
% IOUTB, the two output pins drive 25 ohm termination resistors to AGND,
% and the differential voltage (V_A - V_B) carries the signal.
%
% Compile:  pdflatex 1e_dac_current_mode.tex
% Convert:  pdftocairo -svg 1e_dac_current_mode.pdf
% =============================================================================
\documentclass[border=8pt]{standalone}
%
% =============================================================================
% pmvb-figures.sty -- Poor Man's Validation Bench figure style sheet
% =============================================================================
% Defines the house style for all TikZ figures in PMVB design documents.
%
% Usage in figure source:
%     \documentclass[border=4pt]{standalone}
%     \usepackage{pmvb-figures}
%
% Two color modes:
%     \pmvbDarkMode   -- dark background (default; matches SDD HTML theme)
%     \pmvbLightMode  -- white background (for printed portfolio PDFs)
% =============================================================================
\NeedsTeXFormat{LaTeX2e}
\ProvidesPackage{pmvb-figures}[2026/05/06 PMVB figure styles]
%
% --- Required packages ------------------------------------------------------
\RequirePackage{tikz}
% NOTE: dropped the siunitx option. Our figures use \pmvblabel{} for unit
% annotations rather than circuitikz's SI value formatting, and siunitx.sty
% isn't always available on bare TeX Live installs.
\RequirePackage[american]{circuitikz}
\RequirePackage{xcolor}
\RequirePackage{etoolbox}
%
% --- TikZ libraries ---------------------------------------------------------
\usetikzlibrary{arrows.meta}
\usetikzlibrary{positioning}
\usetikzlibrary{shapes.geometric}
\usetikzlibrary{shapes.misc}
\usetikzlibrary{fit}
\usetikzlibrary{backgrounds}
\usetikzlibrary{calc}
\usetikzlibrary{decorations.pathreplacing}
\usetikzlibrary{patterns}
%
% =============================================================================
% Color palette (FMCW dark theme, matches SDD HTML)
% =============================================================================
\definecolor{pmvbBg}{HTML}{0d1117}
\definecolor{pmvbFg}{HTML}{e6edf3}
\definecolor{pmvbMuted}{HTML}{94a3b8}
\definecolor{pmvbBlue}{HTML}{3b82f6}
\definecolor{pmvbBlueFill}{HTML}{1f2937}
\definecolor{pmvbGreen}{HTML}{10b981}
\definecolor{pmvbGreenFill}{HTML}{1e293b}
\definecolor{pmvbAmber}{HTML}{f59e0b}
\definecolor{pmvbRed}{HTML}{ef4444}
\definecolor{pmvbViolet}{HTML}{a78bfa}
%
\definecolor{pmvbBgLight}{HTML}{ffffff}
\definecolor{pmvbFgLight}{HTML}{0f172a}
\definecolor{pmvbMutedLight}{HTML}{475569}
\definecolor{pmvbBlueFillLight}{HTML}{dbeafe}
\definecolor{pmvbGreenFillLight}{HTML}{d1fae5}
%
% =============================================================================
% Mode selectors
% =============================================================================
\newcommand{\pmvbDarkMode}{%
  \pagecolor{pmvbBg}%
  \color{pmvbFg}%
}
%
\newcommand{\pmvbLightMode}{%
  \colorlet{pmvbBg}{pmvbBgLight}%
  \colorlet{pmvbFg}{pmvbFgLight}%
  \colorlet{pmvbMuted}{pmvbMutedLight}%
  \colorlet{pmvbBlueFill}{pmvbBlueFillLight}%
  \colorlet{pmvbGreenFill}{pmvbGreenFillLight}%
  \pagecolor{pmvbBgLight}%
  \color{pmvbFgLight}%
}
%
\AtBeginDocument{\pmvbDarkMode}
%
% =============================================================================
% Node styles
% -----------------------------------------------------------------------------
% IMPORTANT: each \tikzset{...} block below must be a single argument with NO
% blank lines inside, otherwise \par is inserted and pgfkeys errors out.
% =============================================================================
%
% --- Subsystem (system-level block: a board, a chip, a whole module) ---
\tikzset{
  pmvb subsystem/.style={
    rectangle,
    rounded corners=2pt,
    draw=pmvbBlue,
    line width=0.8pt,
    fill=pmvbBlueFill,
    text=pmvbFg,
    minimum width=28mm,
    minimum height=12mm,
    align=center,
    font=\sffamily\small,
    inner sep=4pt,
  },
}
%
% --- IC / hardware part (functional-level block) ---
\tikzset{
  pmvb ic/.style={
    rectangle,
    rounded corners=1pt,
    draw=pmvbGreen,
    line width=0.6pt,
    fill=pmvbGreenFill,
    text=pmvbFg,
    minimum width=24mm,
    minimum height=10mm,
    align=center,
    font=\sffamily\footnotesize,
    inner sep=3pt,
  },
}
%
% --- Internal block (inside an IC's functional block diagram) ---
\tikzset{
  pmvb internal/.style={
    rectangle,
    draw=pmvbFg,
    line width=0.5pt,
    fill=pmvbBg,
    text=pmvbFg,
    minimum width=18mm,
    minimum height=7mm,
    align=center,
    font=\sffamily\scriptsize,
    inner sep=2pt,
  },
}
%
% --- Pin pad (the box-with-X at IC die boundary in datasheet diagrams) ---
\tikzset{
  pmvb pin/.style={
    rectangle,
    draw=pmvbFg,
    line width=0.4pt,
    fill=pmvbBg,
    minimum size=2.5mm,
    inner sep=0pt,
    path picture={
      \draw[pmvbFg, line width=0.3pt]
        (path picture bounding box.south west) -- (path picture bounding box.north east)
        (path picture bounding box.north west) -- (path picture bounding box.south east);
    },
  },
}
%
% --- Connector (front panel, BNC, banana jack) ---
\tikzset{
  pmvb connector/.style={
    circle,
    draw=pmvbFg,
    line width=0.6pt,
    fill=pmvbBg,
    text=pmvbFg,
    minimum size=8mm,
    font=\sffamily\scriptsize,
    inner sep=1pt,
  },
}
%
% --- Annotation (no border, small muted italic text) ---
\tikzset{
  pmvb annotation/.style={
    text=pmvbMuted,
    font=\sffamily\scriptsize\itshape,
    align=center,
  },
}
%
% --- Subsystem grouping box (used with \node[fit=...]) ---
\tikzset{
  pmvb group/.style={
    rectangle,
    rounded corners=4pt,
    draw=pmvbMuted,
    line width=0.4pt,
    dashed,
    text=pmvbMuted,
    inner sep=8pt,
    font=\sffamily\scriptsize,
  },
}
%
% =============================================================================
% Edge styles
% -----------------------------------------------------------------------------
% Edge styles encode the *kind* of signal flowing on the line.
% =============================================================================
%
\tikzset{
  pmvb digital/.style={
    -Stealth,
    draw=pmvbBlue,
    line width=0.8pt,
    font=\sffamily\scriptsize,
  },
}
%
\tikzset{
  pmvb analog/.style={
    -Stealth,
    draw=pmvbAmber,
    line width=0.8pt,
    font=\sffamily\scriptsize,
  },
}
%
\tikzset{
  pmvb power/.style={
    -Stealth,
    draw=pmvbRed,
    line width=1.0pt,
    font=\sffamily\scriptsize,
  },
}
%
\tikzset{
  pmvb control/.style={
    -Stealth,
    draw=pmvbViolet,
    line width=0.7pt,
    dashed,
    font=\sffamily\scriptsize,
  },
}
%
\tikzset{
  pmvb bus/.style={
    -Stealth,
    draw=pmvbBlue,
    line width=1.4pt,
    double=pmvbBg,
    double distance=0.6pt,
    font=\sffamily\scriptsize,
  },
}
%
% Style for the "xN" bus-width annotation (overlaid on a bus line).
% Dark fill masks the underlying double-line so the white text reads cleanly.
\tikzset{
  pmvb bus width/.style={
    fill=pmvbBg,
    rounded corners=1pt,
    inner sep=2pt,
    font=\sffamily\bfseries\scriptsize,
    text=pmvbFg,
  },
}
%
\tikzset{
  pmvb optional/.style={
    -Stealth,
    draw=pmvbMuted,
    line width=0.6pt,
    dashed,
    font=\sffamily\scriptsize,
  },
}
%
% --- Bidirectional variants ---
\tikzset{
  pmvb digital bi/.style={
    pmvb digital,
    Stealth-Stealth,
  },
}
%
\tikzset{
  pmvb analog bi/.style={
    pmvb analog,
    Stealth-Stealth,
  },
}
%
% =============================================================================
% Diagram-level convenience
% =============================================================================
\tikzset{
  pmvb figure/.style={
    >=Stealth,
    every node/.append style={font=\sffamily},
    node distance=10mm and 14mm,
  },
}
%
% =============================================================================
% Helper macros
% =============================================================================
\newcommand{\pmvblabel}[1]{\textcolor{pmvbMuted}{\sffamily\scriptsize\itshape #1}}
%
\newcommand{\pmvbsubsystem}[3]{%
  node[pmvb subsystem] (#1) {{\bfseries #2}\\\scriptsize #3}%
}
%
\newcommand{\pmvbic}[3]{%
  node[pmvb ic] (#1) {{\bfseries #2}\\\scriptsize #3}%
}
%
% NOTE on bus-width annotations: TikZ's path parser doesn't reliably expand
% \newcommand macros at mid-path positions, so for the "xN" overlay use the
% inline node-in-path syntax directly:
%     \draw[pmvb bus] (A) -- node[pmvb bus width]{$\times 16$} (B);
% The pmvb bus width style is defined above.
%
% -----------------------------------------------------------------------------
% \opamptri{name}{coord}{dir}
%   Draws an op-amp triangle. Apex points right (dir=1) or left (dir=-1).
%   Defines coordinates: <name>_p, <name>_n, <name>_out for wiring.
%   Place any label as a separate \node call after the macro.
%
%   The {coord} argument MUST include its own parentheses, e.g.
%       \opamptri{bufa}{(0.6, -0.5)}{1}
%       \opamptri{tl072}{($(m1e.center) + (4mm, -14mm)$)}{1}
%   This lets callers use either simple coords or calc expressions.
% -----------------------------------------------------------------------------
\newcommand{\opamptri}[3]{%
  \begin{scope}[shift={#2}, xscale=#3]
    \draw[pmvbFg, line width=0.6pt, fill=pmvbBg]
      (0, 0) -- (
%
\ctikzset{bipoles/length=8mm}
%
\begin{document}
\begin{circuitikz}[pmvb figure]

  % ---------------------------------------------------------------------
  % AVDD rail at top
  % ---------------------------------------------------------------------
  \node[font=\sffamily\scriptsize, text=pmvbRed]
    (avdd) at (4, 6) {AVDD = 3.3\,V};

  % ---------------------------------------------------------------------
  % Segmented current source array
  % ---------------------------------------------------------------------
  \node[pmvb ic, minimum width=70mm, minimum height=18mm]
    (csa) at (4, 4.5) {};
  \node[font=\sffamily\bfseries\small, text=pmvbFg]
    at ([yshift=4mm]csa.center) {Segmented PMOS current source array};
  \node[font=\sffamily\scriptsize\itshape, text=pmvbMuted]
    at ([yshift=-1mm]csa.center) {top 9 bits: unary-decoded (one source per code)};
  \node[font=\sffamily\scriptsize\itshape, text=pmvbMuted]
    at ([yshift=-4mm]csa.center) {bottom 3 bits: binary-weighted ${\cdot}$ total = $I_{FS} \approx 20$\,mA};

  \draw[pmvb digital] (avdd.south) -- (csa.north);

  % ---------------------------------------------------------------------
  % 12-bit code switching fabric
  % ---------------------------------------------------------------------
  \node[pmvb ic, minimum width=70mm, minimum height=14mm]
    (sw) at (4, 2.0) {};
  \node[font=\sffamily\bfseries\small, text=pmvbFg]
    at ([yshift=3mm]sw.center) {12-bit code switching fabric};
  \node[font=\sffamily\scriptsize\itshape, text=pmvbMuted]
    at ([yshift=-2mm]sw.center) {each current segment routed to IOUTA or IOUTB by input code};
  \node[font=\sffamily\scriptsize\itshape, text=pmvbMuted]
    at ([yshift=-5mm]sw.center) {(\,$I_A + I_B = I_{FS}$ always)};

  \draw[pmvb digital] (csa.south) -- (sw.north);

  % DB0..DB11 entering from the left
  \node[font=\sffamily\scriptsize, text=pmvbFg, anchor=east, align=right]
    (dbin) at (-2.0, 2.4) {DB0--DB11\\(parallel data from Pico)};
  \draw[pmvb bus] (dbin.east) -- node[pmvb bus width]{$\times 12$} (sw.west |- dbin);

  % CLOCK pin entering from the left, below the bus
  \node[font=\sffamily\scriptsize, text=pmvbFg, anchor=east, align=right]
    (clkin) at (-2.0, 1.5) {CLOCK\\(Pico GP12)};
  \draw[pmvb digital] (clkin.east) -- (sw.west |- clkin);

  % ---------------------------------------------------------------------
  % IOUTA / IOUTB outputs - well separated columns
  % ---------------------------------------------------------------------
  \coordinate (ioutA_pin) at ($(sw.south) + (-2.5, 0)$);
  \coordinate (ioutB_pin) at ($(sw.south) + ( 2.5, 0)$);

  % Drop wires from switch box bottom edges
  \draw[pmvb analog] (ioutA_pin) -- ++(0, -1);
  \draw[pmvb analog] (ioutB_pin) -- ++(0, -1);

  \coordinate (ioutA) at ($(ioutA_pin) + (0, -1)$);
  \coordinate (ioutB) at ($(ioutB_pin) + (0, -1)$);

  % Pin labels (IOUTA / IOUTB)
  \node[font=\sffamily\scriptsize\bfseries, text=pmvbFg, anchor=south, align=center]
    at ($(ioutA_pin) + (0, -0.5)$) {IOUTA};
  \node[font=\sffamily\scriptsize\bfseries, text=pmvbFg, anchor=south, align=center]
    at ($(ioutB_pin) + (0, -0.5)$) {IOUTB};

  % Current direction arrows / labels (to the side)
  \node[font=\sffamily\scriptsize\itshape, text=pmvbAmber, anchor=east]
    at ($(ioutA) + (-2mm, 0)$) {$I_A$};
  \node[font=\sffamily\scriptsize\itshape, text=pmvbAmber, anchor=west]
    at ($(ioutB) + (2mm, 0)$) {$I_B$};

  % Load resistors to ground (well clear of each other)
  \draw (ioutA) to[R=\pmvblabel{25\,$\Omega$}, *-] ($(ioutA) + (0, -2)$) node[ground]{};
  \draw (ioutB) to[R=\pmvblabel{25\,$\Omega$}, *-] ($(ioutB) + (0, -2)$) node[ground]{};

  % V_A / V_B labels on the wire above the resistors (offset to the side)
  \node[font=\sffamily\scriptsize, text=pmvbFg, anchor=east]
    at ($(ioutA) + (-12mm, -1)$) {$V_A = I_A \cdot 25\,\Omega$};
  \node[font=\sffamily\scriptsize, text=pmvbFg, anchor=west]
    at ($(ioutB) + (12mm, -1)$) {$V_B = I_B \cdot 25\,\Omega$};

  % ---------------------------------------------------------------------
  % FSADJ branch on the right side (separate column, no overlap)
  % ---------------------------------------------------------------------
  \coordinate (fsadj_pin)   at ($(csa.east) + (0, 0)$);
  \coordinate (fsadj_corner) at (10.5, 4.5);
  \coordinate (fsadj_R_bot) at (10.5, 2.5);

  \draw[pmvb digital] (fsadj_pin) -- (fsadj_corner);
  \node[font=\sffamily\scriptsize\bfseries, text=pmvbFg, anchor=south]
    at ($(fsadj_pin)!0.6!(fsadj_corner) + (0, 1mm)$) {FSADJ pin};

  \draw (fsadj_corner) to[R=\pmvblabel{$R_{FSADJ}$ = 1.91\,k$\Omega$}, *-]
    (fsadj_R_bot) node[ground]{};

  % ---------------------------------------------------------------------
  % Differential signal + compliance annotation (well below the load Rs)
  % ---------------------------------------------------------------------
  \node[pmvb annotation, align=center, text width=15cm]
    at (4, -2.8)
    {\textbf{Differential signal}\,: $V_A - V_B$ swings $\pm I_{FS} \cdot R_{load} = \pm 0.5$\,V
     ($1\,V_{pp}$) across the input codes.\\[1mm]
     \textbf{Full-scale current}\,: $I_{FS} = 32 \cdot 1.2\,\text{V} / R_{FSADJ}$. With
     $R_{FSADJ}$ = 1.91\,k$\Omega$, $I_{FS} \approx 20$\,mA.\\[1mm]
     \textbf{Compliance window}\,: each output pin must stay within $\pm 1\,$V of ACOM for
     spec'd accuracy. With 20\,mA $\times$ 25\,$\Omega$ = 0.5\,V per leg, we have $\sim$2$\times$ margin.
     Larger $R_{load}$ or $I_{FS}$ would push the PMOS sources out of saturation and
     degrade INL / DNL / THD.};

\end{circuitikz}
\end{document}
